\section{Model, capacity, and main theorem}\label{sec:model}

Let $T$ be an infinite locally finite rooted tree with root $o$ and
levels $L_n=\{v:|v|=n\}$. Assign every edge $e$ a parameter
\begin{equation}\label{eq:channels}
0<\lambda_e\le q<1.
\end{equation}
The tree and the parameters are deterministic and known to the observer.
Start with a uniform root spin in $[4]$ and, independently along edges
conditional on the parent spins, use the transition matrix
\[
P_e(j\mid i)=\lambda_e\mathbf1_{\{j=i\}}+\frac{1-\lambda_e}{4}.
\]
Write $\uvec=(1/4,1/4,1/4,1/4)$, let $p_o^{(n)}$ be the posterior
vector of the root given $\sigma_{L_n}$, and put
\[
a_n=\E\,\TV(p_o^{(n)},\uvec),
\qquad
m_n=4\E\|p_o^{(n)}-\uvec\|_2^2.
\]
Both sequences are nonincreasing by data processing; write their limits as
$a_\infty$ and $m_\infty$.  Reconstruction means $a_\infty>0$, equivalently
$m_\infty>0$.

Put
\begin{equation}\label{eq:rho}
\rho(o)=1,\qquad
\rho(v)=\prod_{e\in[o,v]}\lambda_e^2.
\end{equation}
A nonnegative flow $\theta$ to infinity is conserved at every nonroot
vertex, including vertices with no children. Its mass is
$|\theta|=\sum_{w\in\ch(o)}\theta(o,w)$. For $v\ne o$, write
$e_v$ for its incoming edge. Define
\begin{equation}\label{eq:capacity}
\Cap_3(T,\rho)=\sup\left\{|\theta|:
\sup_{\xi\in\partial T}\sum_{v\in\xi\setminus\{o\}}
\left(\frac{\theta(e_v)}{\rho(v)}\right)^2\le1\right\}.
\end{equation}
The zero flow is allowed. This is the standard $L^3$ convention in which
the path-potential exponent is $3-1=2$.

\begin{theorem}\label{thm:main}
Under \eqref{eq:channels}, for every infinite locally finite deterministic
rooted tree,
\[
\boxed{\quad a_\infty>0\quad\Longleftrightarrow\quad
m_\infty>0\quad\Longleftrightarrow\quad
\Cap_3(T,\rho)>0.\quad}
\]
Put
\[
\delta_q=\frac{(1-q^4)/3000}{(1+9/3000)^3}.
\]
Then the quadratic information satisfies the two-sided comparison
\begin{equation}\label{eq:quantitative-m}
\frac{1-q}{2512}\min\{\Cap_3(T,\rho),1\}
\le m_\infty
\le \frac1{\sqrt{2\delta_q}}\min\{\Cap_3(T,\rho),1\}.
\end{equation}
Consequently,
\begin{equation}\label{eq:quantitative}
\frac{1-q}{20096}\min\{\Cap_3(T,\rho),1\}
\le a_\infty
\le \frac1{2(2\delta_q)^{1/4}}
       \sqrt{\min\{\Cap_3(T,\rho),1\}}.
\end{equation}
\end{theorem}

For $0<\eta\le1$, let $a_n^{(\eta)}$ denote the reconstruction advantage
when every level-$n$ spin is first passed independently through an additional
four-state Potts channel of parameter $\eta$.  Write
$a_\infty^{(\eta)}=\lim_n a_n^{(\eta)}$.

\begin{corollary}[Robust reconstruction]\label{cor:robust}
Under \eqref{eq:channels}, for every fixed $0<\eta\le1$,
\[
 a_\infty^{(\eta)}>0
 \quad\Longleftrightarrow\quad
 \Cap_3(T,\rho)>0.
\]
More precisely,
\begin{equation}\label{eq:robust-lower}
 a_\infty^{(\eta)}
 \ge \frac1{32}
 \min\!\left\{\frac{1-q}{628},\,3\eta^2\right\}
 \min\{\Cap_3(T,\rho),1\}.
\end{equation}
\end{corollary}

The theorem in \cite{OAICapacity} assumes a common channel parameter and a
uniform degree bound. The present proof retains its information-flow
strategy but strengthens both local steps. For necessity, the only imported
probabilistic result is Proposition~3.2 of \cite{OAIThreshold}, stated
precisely in Proposition~\ref{prop:input} below from the source version pinned
in the bibliography. Its source proof uses exact polynomial certificates. We
do not claim to replace that certificate proof. Everything after the stated
input is proved here. In particular, the capacity-to-reconstruction direction
is independent of that input.
